\section{설립 목적과 기능}

한국산업은행은 「한국산업은행법」에 근거하여 설립된 정책금융기관으로, 산업의 개발 및 육성, 사회기반시설 확충, 지역개발, 금융시장 안정, 그리고 지속가능한 성장 촉진을 위한 자금의 공급과 관리를 주요 목적으로 한다. :contentReference[oaicite:0]{index=0}  

이러한 목적은 단순한 금융중개 기능을 넘어, 국가 경제의 구조적 발전을 지원하는 데 초점을 두고 있다는 점에서 일반 상업은행과 구별된다.

\subsection{정책금융기관의 존재 이유}

한국산업은행의 존재는 경제학적으로 \textbf{시장실패(market failure)}를 보완하기 위한 제도적 장치로 이해할 수 있다.

첫째, \textbf{자본시장의 불완전성}이다.  
대규모 산업 투자나 사회기반시설 프로젝트는 초기 투자 규모가 크고 회수 기간이 길기 때문에 민간 금융기관이 적극적으로 참여하기 어렵다. 이러한 영역에서는 자금 공급이 과소하게 이루어지는 경향이 있으며, 이는 경제 전체의 성장 잠재력을 저해할 수 있다.

둘째, \textbf{정보의 비대칭성(asymmetric information)} 문제이다.  
기업의 신용위험이나 사업성에 대한 정보가 불완전한 경우, 금융기관은 보수적으로 대출을 제한하게 된다. 이로 인해 성장 가능성이 높은 기업조차 자금조달에 실패하는 \textit{역선택(adverse selection)} 문제가 발생할 수 있다.

셋째, \textbf{외부효과(externalities)}의 존재이다.  
산업 발전이나 기술 혁신은 개별 기업의 이익을 넘어 사회 전체에 긍정적인 파급효과를 가져오지만, 이러한 효과는 시장에서 충분히 보상되지 않는다. 따라서 민간 부문에만 맡길 경우 투자 수준이 사회적으로 바람직한 수준보다 낮아질 가능성이 크다.

이러한 문제를 해결하기 위해 정부는 정책금융기관을 통해 직접 자금을 공급하거나 금융시장에 개입하게 되며, 한국산업은행은 이러한 역할을 수행하는 핵심 기관이다.

\subsection{핵심 기능}

한국산업은행의 기능은 크게 세 가지로 구분할 수 있다.

\subsubsection{산업금융 공급}

산업은행은 국가 전략 산업 및 기간산업에 대한 자금을 공급함으로써 산업 구조의 형성과 고도화를 지원한다.  
이는 단순한 대출을 넘어, 경제 성장의 방향을 설정하는 역할을 수행한다는 점에서 정책적 의미를 가진다.

\subsubsection{구조조정 및 위기 대응}

경제위기나 기업 부실이 발생할 경우, 산업은행은 구조조정의 중심 기관으로서 기능한다.  
부실기업의 정리, 채권단 조정, 자본 재구성 등을 통해 금융 시스템의 불안정성을 완화하는 역할을 수행한다.

이 기능은 금융시장 안정이라는 공공재(public good)를 제공한다는 점에서 중요하다.

\subsubsection{금융시장 안정화}

산업은행은 시장 불안 시 유동성을 공급하고 금융시장의 충격을 완화하는 역할을 수행한다.  
이는 중앙은행의 통화정책과는 별개로, 보다 미시적이고 산업 중심적인 안정화 기능을 담당한다는 특징이 있다.

\subsection{민간 금융과의 차별성}

한국산업은행은 민간 금융기관과 다음과 같은 차별성을 가진다:

\begin{itemize}
    \item 수익성보다 정책적 목표를 우선시
    \item 장기·대규모 자금 공급 가능
    \item 경제위기 시 적극적 개입 수행
\end{itemize}

그러나 이러한 특성은 동시에 비효율성, 정치적 개입 가능성, 민간 금융 위축(crowding-out) 등의 문제를 야기할 수 있다.

\subsection{소결}

한국산업은행의 설립 목적과 기능은 단순한 금융서비스 제공을 넘어, 시장의 한계를 보완하고 국가 경제의 구조적 발전을 지원하는 데 있다.

따라서 산업은행은 시장 메커니즘과 정부 개입이 결합된 \textbf{혼합적 금융 시스템(hybrid financial system)}의 핵심 구성요소로 이해할 수 있다.